\section{Introduction: Byzantine Consensus in Production Systems}

Byzantine fault tolerance remains a fundamental challenge in distributed systems. The problem, formalized by Lamport et al., asks how a network of processes can reach agreement when some fraction of participants may behave arbitrarily rather than merely crashing. Classical protocols such as PBFT~\cite{castro1999practical} established the theoretical foundation for practical solutions, fixing the $n > 3f$ safety requirement in which $n$ denotes the total node count and $f$ the number of tolerated Byzantine faults. That bound is not an artifact of any particular protocol design: it follows from information-theoretic constraints, since with $n \leq 3f$ Byzantine nodes can partition the network into groups observing contradictory but internally consistent states. The requirement therefore applies to all deterministic Byzantine consensus protocols, independent of timing assumptions or communication pattern.

The decision context for practitioners has shifted, however, from establishing feasibility to selecting among a crowded field of protocol variants. Modern distributed systems demand consensus protocols that scale beyond the limitations of classical approaches. Blockchain networks, distributed databases, and cloud infrastructure all rely on BFT consensus to maintain consistency despite failures, and each imposes distinct throughput, latency, and validator-set-size requirements. Understanding the scalability trade-offs between protocol variants has consequently become essential for practitioners designing production systems, rather than a matter of purely theoretical interest.

The central obstacle is communication cost. Classical PBFT achieves consensus in three phases, pre-prepare, prepare, and commit, and guarantees both safety and liveness in partially synchronous networks. It does so at a cost of $\mathcal{O}(n^2)$ messages per consensus instance, and this quadratic communication complexity becomes a scalability bottleneck as validator sets grow. Protocol design since PBFT has largely been an effort to reduce, amortize, or restructure that cost while preserving the same safety envelope.

The timing model in which a protocol is analyzed materially constrains what can be claimed about it. The synchronous model assumes known message delays, which enables timeout-based protocols but leaves them vulnerable to denial-of-service. The asynchronous model makes no timing assumptions and achieves theoretical robustness, but is limited to probabilistic termination by the FLP impossibility result. The partially synchronous model assumes periods of synchrony after some unknown Global Stabilization Time (GST), and it is this model that underwrites practical protocols with deterministic termination. Comparative claims across protocols are meaningful only when the assumed model is held fixed.

This survey examines advances in Byzantine fault tolerant consensus from 2015 to 2026, with emphasis on scalability approaches, performance characteristics, and formal verification methods. It analyzes how modern variants including Tendermint~\cite{buchman2016tendermint}, HotStuff~\cite{yin2019hotstuff}, and Casper~\cite{buterin2020combining} address the limitations of classical PBFT, and evaluates empirical performance data drawn from production deployments.